\documentclass[11pt,a4paper]{article}

\usepackage[utf8]{inputenc}
\usepackage[T1]{fontenc}
\usepackage{lmodern}
\usepackage{amsmath,amssymb,bm}
\usepackage{graphicx}
\usepackage{booktabs}
\usepackage{hyperref}
\usepackage{geometry}
\usepackage{xcolor}
\usepackage{listings}
\geometry{margin=2.7cm}

\hypersetup{colorlinks=true, linkcolor=blue!50!black, citecolor=blue!50!black, urlcolor=blue!50!black}

\newcommand{\todo}[1]{\par\noindent\textcolor{red!70!black}{\textbf{TODO:} #1}\par}

\title{FYS5419 Project 1: Variational Quantum Eigensolvers and the Lipkin Model}
\author{Your Name}
\date{Spring 2026}

\begin{document}
\maketitle

\begin{abstract}
\todo{Write a concise summary of the problem, methods, main numerical results, and conclusions.}
\end{abstract}

\tableofcontents
\newpage

\section{Introduction}

\todo{Describe the goal of the project: gates, measurements, VQE, entanglement, and the Lipkin model.}

The project studies small quantum systems using both classical diagonalization and variational quantum eigensolver (VQE) methods. The central numerical theme is to compare exact spectra with variational estimates and to interpret changes in eigenstates and entanglement as interaction strengths are varied.

\section{Theory and numerical methods}

\todo{Summarize the variational principle, VQE, Pauli decompositions, density matrices, partial traces, and von Neumann entropy.}

For a normalized trial state $|\psi(\bm\theta)\rangle$, the variational energy is
\begin{equation}
E(\bm\theta)=\langle \psi(\bm\theta)|H|\psi(\bm\theta)\rangle,
\end{equation}
which is an upper bound to the ground-state energy when the trial state belongs to the Hilbert space of the Hamiltonian.

\section{Part a: gates, Bell states, measurements, and entropy}

\todo{Insert implementation details, tables of gate actions, measurement statistics, density matrices, and entropy results.}

\section{Part b: exact solution of the one-qubit Hamiltonian}

\todo{Insert eigenvalue plots, eigenvector-composition plots, and interpretation of the state-character interchange.}

\section{Part c: one-qubit VQE}

\todo{Describe the ansatz, optimizer, expectation-value evaluation, and comparison with exact diagonalization.}

\section{Part d: exact solution of the two-qubit Hamiltonian and entanglement}

\todo{Insert spectra, reduced density matrices, von Neumann entropy, and discussion of level crossings.}

\section{Part e: two-qubit VQE}

\todo{Describe the two-qubit ansatz, optimization results, and comparison with exact diagonalization.}

\section{Part f: Lipkin model and classical diagonalization}

\todo{Show the Pauli decompositions, classical spectra for $J=1$ and $J=2$, and discussion of interaction-strength dependence.}

\section{Part g: Lipkin model with VQE}

\todo{Describe the qubit encoding, ansatz, VQE results, and comparison with exact spectra.}

\section{Discussion}

\todo{Discuss numerical reliability, optimizer convergence, ansatz expressivity, shot noise if used, and physical interpretation.}

\section{Conclusion}

\todo{Summarize what was learned and give concise conclusions.}

\appendix
\section{Use of AI/LLM tools}

\todo{Fill in this appendix if any AI/LLM tools were used. Delete or mark as not applicable only if no such tools were used.}

\subsection{Tools used}

\begin{itemize}
    \item Tool: \todo{tool name, version/access date, access method}
\end{itemize}

\subsection{Text contributions}

\begin{center}
\begin{tabular}{lll}
\toprule
Section & Level & Notes \\
\midrule
Abstract & 0 & -- \\
Introduction & 0 & -- \\
Theory/methods & 0 & -- \\
Results & 0 & -- \\
Discussion/conclusion & 0 & -- \\
\bottomrule
\end{tabular}
\end{center}

\subsection{Code contributions}

\begin{center}
\begin{tabular}{lll}
\toprule
File or notebook & Level & Description \\
\midrule
\texttt{src/fys5419\_project1/...} & 0 & -- \\
\bottomrule
\end{tabular}
\end{center}

\subsection{Verification statement}

All code, equations, and numerical results included in the submitted report have been checked and are understood by the author(s).


\begin{thebibliography}{9}

\bibitem{course-repo}
CompPhysics, \emph{Quantum Computing and Quantum Machine Learning}, course repository.\newline
\url{https://github.com/CompPhysics/QuantumComputingMachineLearning}

\bibitem{hlatshwayo2022}
M.~Q. Hlatshwayo, Y.~Zhang, H.~Wibowo, R.~LaRose, D.~Lacroix, and E.~Litvinova,
``Simulating excited states of the Lipkin model on a quantum computer,''
\emph{Physical Review C} \textbf{106}, 024319 (2022).\newline
\url{https://doi.org/10.1103/PhysRevC.106.024319}

\end{thebibliography}

\end{document}
